% Fichier complété par tools/complete_missing_corrections.py.
% Source : GEO/GEO-03/54_FauteuilRoulant/54_FauteuilRoulant.tex
% Statut : rédigé (3 question(s)).
\begin{corrigebox}[Corrigé rédigé]
\CorrigeQuestion{1}
On identifie les solides, les pivots des roues et le contact de
                roulement avec le sol. Le graphe est bâti/châssis--roues par pivots, puis
                roues--sol par contacts de roulement.
\CorrigeQuestion{2}
Pour chaque roue, la condition de roulement donne
                $v=R\omega$. Pour deux roues motrices séparées de la voie $2a$, la vitesse
                du centre et la vitesse de lacet sont
                $v_G=R(\omega_d+\omega_g)/2$ et
                $\dot\psi=R(\omega_d-\omega_g)/(2a)$.
\CorrigeQuestion{3}
La trajectoire est rectiligne si $\omega_d=\omega_g$, une rotation sur
                place si $\omega_d=-\omega_g$, et un cercle de rayon
                $\rho=a(\omega_d+\omega_g)/(\omega_d-\omega_g)$ sinon.
\end{corrigebox}
